% No preamble. Prefix pid-. Canonical normalized periodization only.
\section{Bounded finite periodization of analytic interaction classes}
\label{pid-section}

Use the infinite activity space and normalized-total quotient of
Notebook29. Fix the cell size $J$, at most $D$ scalar coordinates
per cell, exponent $a>0$, and radius $b>0$. Carriers are nonempty
finite connected subsets of the cell graph $\mathbb Z^d$, and
\begin{equation}
 \|E\|_{a,b}=\sup_i\sum_{H\ni i}e^{a|H|}\|E_H\|_b,
 \qquad
 \Phi_E(A)=\sum_H[E_H(A|_H)-E_H(0)]
 \label{pid-space}
\end{equation}
on compact fields. Let $\Lambda_N=(\mathbb Z/N\mathbb Z)^d$;
the underlying current lattice has period $NJ$. Write $\pi_N$
for cell and coordinate projection and $A^{\rm per}$ for the
infinite periodic lift of a finite field. It is bounded, generally
not summable. Internal link directions and color types are
identified by their actual periodic conventions.

\begin{theorem}[Canonical normalized periodization]
\label{pid-theorem}
Every class with a $J$-cell-translation-invariant normalized total
has a translation-covariant representative without increasing the
norm of any chosen input representative. There is a canonical
linear periodization map
\begin{equation}
 \operatorname{Per}_N:\mathcal Q^{\rm tr}_{a,b}
                   \longrightarrow\mathcal Q^N_{a,b},
 \qquad \|\operatorname{Per}_N\|\le1,
 \label{pid-operator}
\end{equation}
where the finite quotient uses connected torus carriers and
equality of vacuum-normalized totals. Its scalar functional is
holomorphic on the original open radius-$b$ polydisc. Every
fixed derivative tensor is the exact relative periodization of
the infinite total tensor. Total spatial and color symmetries
pass whenever they preserve the torus. Local gauge invariance
passes on the specified common logarithmic chart.
\end{theorem}

The proof retains raw constants in carrier norms and normalizes
only the total. It does not identify the result with an unrelated
actual finite-regulator quantum output.

\subsection{Translation averaging without norm loss}
\label{pid-averaging-section}

Let $\tau_t$ translate carriers, their coordinate slots and their
functions by a cell-lattice vector $t$. It is an isometry of the
activity space. If $[E]$ is translation invariant, then
$[\tau_tE]=[E]$. For centered F\o lner cubes $F_m$, put
\begin{equation}
 E^{(m)}=\frac1{|F_m|}\sum_{t\in F_m}\tau_tE.
 \label{pid-Cesaro}
\end{equation}
Its norm is at most $\|E\|_{a,b}$, and for every fixed $s$,
\begin{equation}
 \|\tau_sE^{(m)}-E^{(m)}\|_{a,b}
 \le\frac{|(F_m+s)\mathbin\triangle F_m|}{|F_m|}
                                      \|E\|_{a,b}\longrightarrow0.
 \label{pid-Folner-bound}
\end{equation}
Each fixed carrier has finitely many input coordinates. Countable
Montel extraction gives locally uniform limits $E_H^{\rm cov}$
for all carriers. Fatou, first on finite carrier collections and
smaller polydiscs and then at the original open radius, yields
\begin{equation}
 \|E^{\rm cov}\|_{a,b}\le\|E\|_{a,b}.
 \label{pid-covariant-bound}
\end{equation}
Equation~\eqref{pid-Folner-bound} passes componentwise and gives
exact covariance. On a compact input supported in $K$, the
omitted carriers with $|H|>Q$ contribute at most
$2|K|\|E\|_{a,b}e^{-aQ}$. The finite remaining sum converges,
so $\Phi_{E^{\rm cov}}=\Phi_E$. No uniqueness of this covariant
representative is needed.

Only cell translations are averaged. Finer unit translations,
point symmetries and color invariance need not hold for the
individual representative and will pass through its total.

\subsection{The exact orbit count and wrapped carrier norm}
\label{pid-orbit-section}

For a covariant representative define
\begin{equation}
 U_N^E(A)=\sum_{i\in\Lambda_N}\sum_{H\ni i}
                 \frac{E_H(A^{\rm per}|_H)}{|H|},
 \qquad
 \Phi_N^{[E]}(A)=U_N^E(A)-U_N^E(0).
 \label{pid-density}
\end{equation}
Choose one fundamental representative for each cell $i$. The
formula is independent of that choice by covariance. Its series
is absolutely convergent in the supremum norm on the full open
$b$ polydisc, bounded by $N^d\|E\|_{a,b}$ before normalization.

Let $\mathcal O$ be an orbit of finite carriers under
$N\mathbb Z^d$. A nonempty finite carrier has no nonzero
translation stabilizer. Every point in a representative $H$
corresponds to exactly one incidence pair $(i,H')$ with
$i\in\Lambda_N$, $H'\in\mathcal O$ and $i\in H'$. Hence
\begin{equation}
 \sum_{i\in\Lambda_N}\sum_{\substack{H'\in\mathcal O\\ i\in H'}}
                      \frac1{|H'|}=1.
 \label{pid-incidence}
\end{equation}
Covariance and periodicity make the summand value constant on
the orbit. Therefore
\begin{equation}
 U_N^E(A)=\sum_{\mathcal O}E_H(A^{\rm per}|_H),
                   \qquad H\text{ any representative of }\mathcal O.
 \label{pid-orbit-sum}
\end{equation}
This counts each carrier orbit once, including those whose
projection is not injective.

For a connected nonempty torus carrier $K$, set
\begin{equation}
 E_{N,K}(A|_K)=
       \sum_{\mathcal O:\,\pi_NH=K}E_H(A^{\rm per}|_H).
 \label{pid-grouped-family}
\end{equation}
Coordinate repetition is a linear map of maximum norm one.
Every summand is holomorphic on the original $b$ polydisc and
strictly supported on $K$. Projection preserves connectedness.

Fix a torus cell $j$ and one infinite lift. For every orbit
meeting $j$, choose one representative containing that lift.
Distinct orbits give distinct carriers. Since
$|\pi_NH|\le|H|$, this gives
\begin{align}
 \sum_{K\ni j}e^{a|K|}\|E_{N,K}\|_b
 &\le\sum_{\mathcal O:\,j\in\pi_NH}
                         e^{a|\pi_NH|}\|E_H\|_b\notag\\
 &\le\sum_{H\ni j}e^{a|H|}\|E_H\|_b
 \le\|E\|_{a,b}.
 \label{pid-wrapped-bound}
\end{align}
This proves normal convergence of the grouped family and the
no-loss carrier estimate. Choosing one incidence here is an
upper bound for the norm; the exact energy multiplicities were
already counted in \eqref{pid-incidence}. No factor
$|\pi_NH|/|H|$ belongs in the energy.

The raw grouped coefficients retain their constants. Normalizing
each coefficient separately could cost a factor two, and is
unnecessary. For example the constant singleton family
$E_{\{i\}}=c$ is the zero normalized class, but its raw periodic
total is $N^dc$. A canonical unnormalized exponential therefore
needs an additional vacuum convention.

\subsection{Derivatives, representative independence and local stability}
\label{pid-kernel-section}

Let $K_n^E(i_1,\alpha_1;\ldots;i_n,\alpha_n)$ be the infinite
total derivative kernel, with cell indices and internal types
separated. At every fixed $n$, connected-carrier Cauchy estimates
give absolute anchored summability with a positive spatial
exponent. Constants can depend on $n,D,J,a,b$; no uniform
all-order absolute Taylor norm is being asserted.

Differentiate \eqref{pid-orbit-sum}. A derivative in a periodic
coordinate differentiates each old occurrence of that coordinate.
Translate its first occurrence $i_1+Nz_1$ to a fixed lift $i_1$.
Every carrier through $i_1$ occurs exactly once. The remaining
slots then range over all their periodic lifts. Thus
\begin{align}
 &D^nU_N^E(0)[e_{i_1,\alpha_1},\ldots,e_{i_n,\alpha_n}]
 \notag\\
 &\qquad=\sum_{z_2,\ldots,z_n\in\mathbb Z^d}
 K_n^E(i_1,\alpha_1;
       i_2+Nz_2,\alpha_2;\ldots;i_n+Nz_n,\alpha_n).
 \label{pid-kernel-periodization}
\end{align}
Repeated finite slots include both coincident and distinct
infinite lifts with their ordinary multilinear multiplicities.
Absolute anchored summability justifies the rearrangement.

A null class has all $K_n=0$ for $n\ge1$. Its normalized
periodic total has every Taylor coefficient zero and hence is
zero on the connected finite-dimensional $b$ polydisc. Two
covariant representatives of the same class therefore give the
same normalized energy. This proves independence of the
translation-averaging subsequence. The scalar map is linear
even if representative selection is not. Taking representatives
approaching the quotient infimum in
\eqref{pid-covariant-bound},\eqref{pid-wrapped-bound} proves
\eqref{pid-operator}, and
\begin{equation}
 |\Phi_N^{[E]}(A)|\le2N^d\|[E]\|_{\mathcal Q}.
 \label{pid-energy-bound}
\end{equation}

There is exact scalar tiling consistency. For an integer $q>0$,
repeat $A$ on $\Lambda_{qN}$. The density in
\eqref{pid-density} is then $N$-periodic, so
\begin{equation}
 \Phi_{qN}^{[E]}(A\text{ repeated})=q^d\Phi_N^{[E]}(A).
 \label{pid-tiling}
\end{equation}
A rectangular multiple has its corresponding volume factor.

There is also eventual exact stabilization of each fixed lifted
coefficient. Let $K$ be a fixed finite connected infinite carrier
and choose, conservatively, $N>4\operatorname{diam}(K)+4$.
The copies $K+Nz$ are disjoint and nonadjacent. A connected $H$
with $\pi_NH=\pi_NK$ lies in their union, hence in one copy;
surjectivity of its projection makes it equal to that copy.
Equation~\eqref{pid-grouped-family} consequently has one orbit,
and under lifted coordinates
\begin{equation}
 E_{N,\pi_NK}=E_K
 \label{pid-fixed-carrier}
\end{equation}
on the full original polydisc. This gives local normal convergence
of canonical periodizations without a winding-tail estimate for
fixed $K$. It asserts neither global norm convergence nor
stabilization on carriers growing with $N$.

\subsection{Local gradients and exact finite gauge invariance}
\label{pid-gauge-section}

For a fixed infinite scalar coordinate $i$ and any bounded field
$\|A\|_\infty\le\rho<b$, define
\begin{equation}
 G_i^E(A)=\sum_{H\ni\operatorname{cell}(i)}
                                 \partial_iE_H(A|_H).
 \label{pid-gradient}
\end{equation}
Directional Cauchy gives
\begin{equation}
 |G_i^E(A)|\le\frac{\|E\|_{a,b}}{b-\rho}.
 \label{pid-gradient-bound}
\end{equation}
If $A,B$ agree in a $Q$-cell neighborhood of $i$, each differing
summand has a connected carrier of size at least $cQ$. Thus
\begin{equation}
 |G_i^E(A)-G_i^E(B)|
       \le\frac{2\|E\|_{a,b}}{b-\rho}e^{-caQ}.
 \label{pid-gradient-tail}
\end{equation}
This proves continuity under compact truncation of bounded
fields. On compact fields $G_i^E$ is the derivative of
$\Phi_E$. A null family therefore has $G_i^E=0$ on all
bounded interior fields, even without covariance of that
representative.

For a covariant representative, the same first-slot reanchoring
used above works at any periodic interior field and gives
\begin{equation}
 \partial_iU_N^E(A)=G_i^E(A^{\rm per}).
 \label{pid-periodic-gradient}
\end{equation}
This is another proof that null classes periodize to constant
raw totals and to zero normalized totals.

Suppose the infinite normalized total is locally gauge invariant.
For a fixed unit-current site $v$, its infinitesimal logarithmic
gauge field $\mathcal L_v(A)$ reads only the incident bond
potentials and is supported on those finitely many bonds.
On compact fields the Ward identity is
\begin{equation}
 \sum_{i\text{ incident to }v}
                    G_i^E(A)\mathcal L_{v,i}(A)=0.
 \label{pid-local-Ward}
\end{equation}
Compactly truncate an arbitrary bounded interior $A$. The local
generator values stabilize, and \eqref{pid-gradient-tail}
passes the gradient terms. Thus \eqref{pid-local-Ward} holds
for bounded periodic fields too. Using
\eqref{pid-periodic-gradient}, identify all incident orientations
and periodic multiplicities. The finite total obeys
\begin{equation}
 D\Phi_N^{[E]}(A)[\mathcal L_{\lambda,N}(A)]=0
 \label{pid-finite-Ward}
\end{equation}
for every finite torus gauge parameter. No noncompact infinite
gauge path has been applied to $\Phi_E$; only a local gradient
identity was extended by absolute convergence.

The logarithmic generator must be holomorphic on the domain.
If $b$ lies inside its fixed chart radius, the assertion holds
on the whole open $b$ polydisc. Otherwise keep full-$b$
holomorphy of the interaction and restrict the gauge assertion
explicitly to $b_g=\min(b/2,r_{\log}/2)$, where $r_{\log}$ is
a fixed radius of holomorphy for that generator. The actual
small-field setup already chooses $b$ inside this chart. Finite
gauge invariance follows along any gauge path whose fields stay
in the common domain. Large paths leaving it are not covered.

All finite nonlinear Ward equations also pass coefficientwise
in field degree. At any fixed order, differentiating the compact
infinite identity uses finitely many local generator coefficients,
nearest-neighbor shifts and coincident-slot embeddings of total
kernels. Equation~\eqref{pid-kernel-periodization} commutes with
these operations: summing images turns an infinite coincidence
condition into the exact periodic one. Fixed-order summability
justifies it. Equivalently, differentiate the exact analytic
identity \eqref{pid-finite-Ward}.

\subsection{Other symmetries and the finite-history boundary}
\label{pid-scope-section}

Any finer translation, hypercubic or global color symmetry of
the infinite total is inherited by its kernels. Relative
periodization commutes with each symmetry preserving the torus.
Taylor uniqueness gives the same symmetry of the normalized
finite total. Negative bonds use the signed reversal of their
positive-bond logarithmic coordinate. No symmetry of individual
$E_{N,K}$ is needed. Reality follows similarly; a real covariant
representative can also be obtained by averaging with its
complex-conjugate reflection without norm increase.

This completes Theorem~\ref{pid-theorem}. The canonical finite
functionals have exact Ward identities and a no-loss connected
representation. They need not equal an unrelated actual
finite-regulator family. Two uniformly bounded finite sequences
can have the same infinite local interaction while differing
by winding-supported terms at every finite size. Compactness
does not assert those terms are exactly absent.

Consequently, inserting $\operatorname{Per}_N[E]$ into a quantum
integral requires either an actual history/normalization comparison
or a construction using these periodized inputs from the outset.
Stationary backgrounds, retained windows, Gaussian normalizations
and gauge factors are not identified by periodizing the interaction
alone. Raw exponentials also require a separately specified vacuum
convention. No quantum integration identity, continuity of the
actual quantum update, or infinite induction is proved here.
